\chapter*{Wstęp}\addcontentsline{toc}{chapter}{Wstęp}\markboth{Wstęp}{Wstęp}

Ta książka zaczyna się od arytmetyki, której nie da się obejść. Doba ma dwadzieścia cztery godziny. Jeśli dwanaście z nich zajmuje sen, zostaje dwanaście, i wszystko, co chcesz zrobić w ciągu dnia, musi się w nich zmieścić: praca, dojazd, jedzenie, dom, ruch, rozmowy i chleb. Nie ma tu żadnej sztuczki. Całe planowanie polega na tym, żeby ta reszta była policzona, a nie wyczuwana.

Tytuł brzmi jak obietnica, ale tekst jest raczej zbiorem pytań. Ile snu potrzebujesz naprawdę? Czy dwanaście godzin to Twoja potrzeba, czy skutek długiego zmęczenia? Co w Twoim tygodniu zjada najwięcej czasu i czy da się to zmienić? Gdzie w dobie mieści się chleb, który wymaga czasu biernego, a ten przypada na noc?

\section*{Dla kogo jest ta książka}

Piszę dla dorosłego, który lubi liczby i woli zobaczyć rachunek niż slogan. Mam na myśli trzy rodzaje czytelników. Pierwszy śpi za mało od lat, wie o tym i chce zaplanować sen zamiast odrabiać go w weekendy. Drugi wie od lekarza albo z własnego doświadczenia, że jego potrzeba snu jest długa, i nie zamierza z nią walczyć. Trzeci piecze chleb, zauważył, że dobry chleb wymaga czasu, którego nie ma w czasie aktywnym, i chce ten czas bierny wykorzystać w dobie.

Dwanaście godzin to dużo. Typowe zalecenia dla dorosłych mieszczą się wyraźnie poniżej tej wartości: zwykle podaje się około \verify{7--9 godzin snu na dobę}. Nie piszę tego, żeby kogokolwiek pocieszać ani straszyć. Dwanaście godzin to założenie planu, a nie norma, którą każdy powinien przyjąć. Jeśli Twoja potrzeba jest inna, cała arytmetyka przesunie się o tyle, o ile ją zmienisz. Książka jest napisana tak, żeby ją przeliczyć.

\section*{Premisa}

Plan opiera się na dwóch założeniach. Pierwsze: sen jest stały i traktowany jak zobowiązanie, a nie jak resztka po wieczorze. Drugie: chleb pasuje do takiego rytmu lepiej, niż się wydaje. Ciasto na zakwasie albo na drożdżach może długo fermentować w lodówce, nawet przez kilkanaście godzin. Czas, w którym ciasto rośnie, jest czasem biernym: nie wymaga stania przy blacie, więc może przypaść na sen. W planie zakładamy, że czas aktywny przy jednym bochenku to około kwadransa rozłożonego na kilka momentów dnia.

To jest argument o organizacji czasu, a nie o zdrowiu. Nie twierdzę, że chleb poprawia sen ani że sen poprawia chleb. Twierdzę tylko, że oba zajęcia dadzą się w jednej dobie ułożyć, jeśli policzysz każdą godzinę.

To nie jest porada medyczna. Nie diagnozuję i nie leczę. Jeśli śpisz długo, a mimo to jesteś zmęczony, jeśli chrapiesz, budzisz się z bólem głowy albo sen nie daje Ci odpoczynku, zobacz lekarza. Długi sen \verify{bywa objawem problemów zdrowotnych}, a plan go nie wyjaśni.

\section*{Co zawiera książka}

Książka ma \val{chapters.total} rozdziałów w trzech częściach. Część pierwsza, \emph{Arytmetyka doby}, pokazuje, jak policzyć dobę, ile snu potrzebujesz i co realnie zjada czas. Część druga, \emph{Dzień wokół snu}, dotyczy pracy, rytmu snu, szablonu dnia i domu. Część trzecia, \emph{Chleb}, wyjaśnia fermentację na zimno i zawiera harmonogramy wypieków. Na końcu są dodatki: receptury z procentami piekarskimi, szablony doby do wypełnienia, źródła i weryfikacja oraz rejestr stanu książki.

Każdy rozdział ma cele, rozpisany przykład z liczbami, jedną pułapkę, do której wpada wiele osób, ćwiczenia z odpowiedziami w Dodatku A i jedno zadanie na tydzień. Liczby, które książka drukuje, są liczone z danych, a nie wpisywane ręcznie. Twierdzenia o świecie, które wymagają źródła, oznaczono znaczkiem \verify. Ich liczbę i sposób sprawdzania opisują rejestr i Dodatek D.

\section*{Czego ta książka nie obiecuje}

Nie obiecuję, że po tej książce będziesz spać lepiej, pracować wydajniej, schudniesz ani że Twój chleb wyjdzie za pierwszym razem. Nie leczę bezsenności i nie zastępuję diagnozy. Nie podaję godzin, w których zrobisz konkretne rzeczy, bo takie tabele są przybliżone i zależą od tego, kto je wypełnia. Nie wiem, jak plan zadziała u Ciebie, a książka nie została jeszcze przejrzana przez lekarza medycyny snu. Ten stan jest zapisany w nagłówku każdej strony.

Koszt jest realny. Dwanaście godzin snu to połowa doby. Każda godzina, którą śpisz, to godzina, której nie masz na coś innego, i ta książka każe zapłacić za to własnym czasem: częścią wieczorów, części zakupów i części zajęć, które lubisz. Jeśli po lekturze uznasz, że cena jest za wysoka, to też jest wynik, i warto go zapisać.

Jeśli nie masz pewności, czy dwanaście godzin to dla Ciebie właściwa liczba, zanim zaplanujesz tydzień, porozmawiaj z lekarzem. Dodatek D opisuje, jak książka traktuje tę kwestię i dlaczego nie chce być w niej ostatnim słowem.
